\section{Day 4: Conviction, Repentance, and Remission}

\dayaxis{John 16:7--15; John 20:19--23; Acts 2:37--41}{Truthful contrition and confidence in Christ's forgiveness}{Examine conscience; arrange sacramental confession if grave sin or prudent spiritual need requires it.}

\subsection{Read: John 16:7--15 and Acts 2:37--41}

The promised Advocate exposes the world's false judgment of sin, righteousness, and condemnation. This conviction is not the voice of hopeless accusation. At Pentecost Peter's hearers are “cut to the heart,” ask what they must do, and receive an answer: repent, be baptized, receive forgiveness and the Holy Spirit, enter the promise. The wound of truth opens toward sacramental mercy.

John's Gospel joins the risen Christ's breathing of the Spirit to the apostolic ministry of forgiving and retaining sins (Jn 20:19--23). Pentecostal devotion is therefore malformed if it seeks power while evading repentance, or claims private absolution apart from the means Christ entrusted to the Church. For a baptized Catholic conscious of mortal sin, perfect contrition includes the intention of sacramental confession as soon as reasonably possible; ordinarily, sacramental absolution restores the lost state of grace.

Conviction also differs from scrupulosity. The Spirit names sin in order to heal and reorder; scrupulosity often multiplies doubtful guilt, distrusts absolution, and makes attention circle the self. A confessor and qualified clinician can help where obsessive anxiety is present. The novena should not intensify it.

\subsection{Meditation: fire that purifies rather than annihilates}

Isaiah's burning coal touches and cleanses the prophet before mission (Isa 6:1--8). Pentecostal fire likewise signifies holy presence and purifying charity, not divine delight in pain. God does not need suffering in order to become willing to forgive; the Son has offered the perfect sacrifice, and the Spirit applies its fruit.

Make the examination concrete. Name actions and omissions, not a vague sense of worthlessness. Ask where knowledge, freedom, matter, habit, fear, or coercion affects culpability without using diminished culpability to rename evil good. Then ask what repair love requires: confession, apology, repayment, correction of a false report, removal of an occasion of sin, professional help, or patient amendment.

Pray also for those who cannot yet recognize a wound they caused. Intercession is not permission to control their conversion. Entrust them to truth and mercy, maintain necessary boundaries, and refuse both vengeance and denial.

\novenaexamination{What sin do I excuse because it serves my preferred identity? What forgiven sin do I keep resurrecting because I distrust mercy? What concrete repair belongs to repentance today?}

\bilingualprayer{Proper collect for Day 4}{Project-composed Latin prayer and project English translation; not an approved liturgical collect.}{%
Pater misericordiárum, qui per Fílium tuum Spíritum remissiónis effudísti: illúmina consciéntias nostras, da nobis contritiónem veram, et ad fontem reconciliatiónis nos perdúc; ut, peccátis dimíssis, in novitáte vitæ ambulémus et misericórdiam quam accépimus áliis largiámur. Per Christum Dóminum nostrum. Amen.}{%
Father of mercies, who through thy Son hast poured forth the Spirit of forgiveness: enlighten our consciences, give us true contrition, and lead us to the fountain of reconciliation; that, our sins forgiven, we may walk in newness of life and extend to others the mercy we have received. Through Christ our Lord. Amen.}


\prayerinstruction{Continue with \textit{Veni Creator Spiritus}, the Lord's Prayer, and the Lesser Doxology.}
